% Hafta 5 — Yönetilen dil neyi çözer, neyi çözmez?
% Derleme: py -3.12 tools/latex/derle.py docs/week-5/assets/h05-02-yonetilen-dil.tex
\documentclass[tikz,border=8pt]{standalone}
\usepackage{cen429-cizim}
\begin{document}
\begin{tikzpicture}
  \node[baslik] (b) at (0,0) {Yönetilen dil neyi çözer, neyi çözmez?};
  \node[altbaslik] at (b.south west) {Java/Kotlin/C\# çöp toplayıcı + sınır denetimi};
  \node[iyikutu=70mm, anchor=north west] (sol) at (0,-2.0)
    {\kb{ÇÖZER (dil düzeyinde)}\\[2mm]
     tampon taşması\\use-after-free\\çift serbest bırakma\\ilklendirilmemiş değişken\\[3mm]
     $\rightarrow$ bellek güvenliği};
  \node[kotukutu=70mm, right=14mm of sol] (sag)
    {\kb{ÇÖZMEZ}\\[2mm]
     enjeksiyon (SQL, komut, XML)\\güvensiz deserialization\\koda gömülü sırlar\\zafiyetli bağımlılıklar\\mantık/yetki hataları};
  \node[fit=(sol)(sag), inner sep=0pt] (grp) {};
  \node[dip, text width=155mm, anchor=north west] at ($(grp.south west)+(0,-8mm)$)
    {Dil bir hata SINIFINI kaldırır; kalan sınıflar için hâlâ siz sorumlusunuz.};
\end{tikzpicture}
\end{document}
